\documentclass[12pt,a4paper]{report}

\usepackage[utf8]{inputenc}
\usepackage[T1]{fontenc}
\usepackage[french]{babel}
\usepackage{graphicx}
\usepackage{hyperref}
\usepackage{listings}
\usepackage{xcolor}
\usepackage{geometry}
\usepackage{titlesec}
\usepackage{fancyhdr}
\usepackage{tcolorbox}
\usepackage{parskip}
\usepackage{setspace}

\geometry{a4paper, margin=2.5cm}
\setstretch{1.2}

\definecolor{codegreen}{rgb}{0,0.6,0}
\definecolor{codegray}{rgb}{0.5,0.5,0.5}
\definecolor{codepurple}{rgb}{0.58,0,0.82}
\definecolor{backcolour}{rgb}{0.95,0.95,0.92}

\lstdefinestyle{mystyle}{
    backgroundcolor=\color{backcolour},   
    commentstyle=\color{codegreen},
    keywordstyle=\color{magenta},
    numberstyle=\tiny\color{codegray},
    stringstyle=\color{codepurple},
    basicstyle=\ttfamily\footnotesize,
    breakatwhitespace=false,         
    breaklines=true,                 
    captionpos=b,                    
    keepspaces=true,                 
    numbers=left,                    
    numbersep=5pt,                  
    showspaces=false,                
    showstringspaces=false,
    showtabs=false,                  
    tabsize=2
}
\lstset{style=mystyle}

\begin{document}

\begin{titlepage}
    \begin{center}
        \vspace*{1cm}
        
        \Huge
        \textbf{Rapport de Projet : Intelligence Artificielle}
        
        \vspace{0.5cm}
        \LARGE
        \textbf{Cellule SMA - Système Multi-Agents pour la Veille Stratégique}
        
        \vspace{1.5cm}
        
        \textbf{Automatisation intelligente via LangGraph et RAG Hybride}
        
        \vfill
        
        \Large
        \textbf{Réalisé par :}\\
        CHOUKHAIRI Noureddine\\
        ELATTAR Abderahmane\\
        MOUIH Anass\\
        OUEDGHIRI BEN OTMANE Abdelmalek\\
        
        \vspace{1.5cm}
        
        \textbf{École Marocaine des Sciences de l'Ingénieur (EMSI)}\\
        \vspace{0.5cm}
        Date : Juin 2026
        
    \end{center}
\end{titlepage}

\tableofcontents
\newpage

\chapter{Introduction et Problématique}

\section{Contexte de la Veille Stratégique}
Dans un environnement économique et technologique en perpétuelle mutation, la veille stratégique est devenue une nécessité vitale pour la survie et la compétitivité des entreprises. Le contexte actuel impose aux décideurs de traiter un volume croissant d'informations tout en garantissant leur absolue fiabilité. L'enjeu majeur est de transformer cette masse de données hétérogènes en intelligence actionnable.

La veille manuelle (qu'elle soit technologique, concurrentielle ou réglementaire) souffre de plusieurs maux fondamentaux. Elle est lente, coûteuse et extrêmement chronophage. Ce processus manuel réduit drastiquement la capacité d'une entreprise à réagir rapidement aux signaux faibles de son marché. 

\section{La Fragmentation des Données}
Une difficulté supplémentaire réside dans la fragmentation des sources d'information. Les données stratégiques d'une entreprise sont typiquement séparées en deux mondes étanches :
\begin{itemize}
    \item \textbf{Le monde interne :} Des documents confidentiels (PDFs de roadmap, bilans financiers, mémos de gouvernance) stockés sur des serveurs sécurisés.
    \item \textbf{Le monde externe :} Le web en temps réel (articles scientifiques, tendances open-source, flux d'actualités mondiales).
\end{itemize}

L'enjeu décisionnel est complexe : comment croiser ces deux mondes rapidement pour prendre des décisions stratégiques (comme l'allocation d'un budget ou la modification d'une roadmap) sans compromettre la sécurité des données internes ?

\section{Objectif du Projet}
Face à ces défis, l'objectif de notre projet est de concevoir et de développer une "Cellule de Veille Stratégique Automatisée". Il s'agit d'un système intelligent s'appuyant sur l'intelligence artificielle générative pour agréger, analyser et synthétiser de l'information en un temps record. 

Cependant, comme nous le verrons dans le chapitre suivant, l'utilisation brute de l'Intelligence Artificielle générative pose des risques majeurs que ce projet s'est fait un devoir de résoudre grâce à une architecture Multi-Agents (SMA) avancée.

\chapter{Limites de l'IA Classique et l'Approche Hybride}

\section{Le Piège du Chatbot Simple}
L'essor des Large Language Models (LLM) comme ChatGPT a laissé penser que la veille pouvait être entièrement déléguée à un simple prompt. Cependant, un chatbot classique ne peut garantir une veille rigoureuse à lui seul. Il répond de manière générique, il ne possède pas le contexte métier de l'entreprise (ses documents confidentiels) et, surtout, ses connaissances sont figées dans le temps, ce qui le rend incapable de réagir à l'actualité en temps réel.

\section{Le Risque Majeur : L'Hallucination}
L'obstacle le plus critique à l'adoption des LLMs dans la prise de décision d'entreprise est le phénomène d'hallucination. Les LLMs ont tendance à inventer des faits, des chiffres et des citations lorsqu'ils manquent de données ou lorsqu'ils tentent de combler un vide statistique dans leurs probabilités. 

Dans un contexte de veille stratégique, ce comportement représente un risque majeur. Aucune marge d'erreur n'est acceptable lorsqu'il s'agit de voter des budgets ou de modifier des roadmaps. L'enjeu critique du système est de garantir une information $100\%$ vérifiable et sourcée.

\section{La Solution : Le RAG Hybride}
Pour contrer ces limites, notre solution s'appuie sur le concept de RAG (Retrieval-Augmented Generation) que nous avons adapté à l'entreprise sous la forme d'un \textbf{RAG Hybride}. Le but est simple : empêcher l'IA de réfléchir "dans le vide". Nous la forçons à synthétiser des textes qui existent réellement et que nous lui fournissons explicitement.

Le RAG Hybride se divise en deux parties :
\begin{enumerate}
    \item \textbf{Le RAG Interne (La Mémoire) :} Une base vectorielle persistante (ChromaDB) qui stocke les documents internes confidentiels de l'entreprise.
    \item \textbf{Le RAG Externe (Les Yeux) :} Des outils connectés au web en temps réel (APIs GitHub, arXiv, flux RSS) qui captent l'actualité à la seconde près.
\end{enumerate}

\chapter{Architecture Multi-Agents et LangGraph}

\section{Pourquoi LangGraph face à LangChain Classique ?}
La majorité des applications basées sur l'IA utilisent LangChain pour exécuter des actions linéairement. Cependant, une architecture linéaire ne permet pas de gérer l'incertitude ni de s'auto-corriger. Si une erreur survient en bout de chaîne, le système entier échoue et délivre une mauvaise réponse.

\textbf{LangGraph} résout ce problème en introduisant la notion de graphes cycliques. Il permet de créer des flux de travail incluant des boucles de rétroaction (feedback loops). Cette capacité de retour à des étapes antérieures est idéale pour l'auto-correction et la validation itérative, ce qui est le cœur de notre système anti-hallucination.

\section{Le Flux de Données (DFD)}
Le workflow de notre système se découpe en 4 grandes phases :
\begin{itemize}
    \item \textbf{Orchestration :} Pilotage et routage intelligent dès le lancement de la requête.
    \item \textbf{Fan-out :} Exécution asynchrone et parallèle de la recherche interne et externe pour diviser le temps de traitement.
    \item \textbf{Fan-in :} Analyse comparative stratégique (convergence des données).
    \item \textbf{Validation :} Contrôle anti-hallucination avec boucle de correction conditionnelle.
\end{itemize}

\section{Les Agents de Recherche}
Le cœur du système de veille repose sur trois agents spécialisés dans la collecte et l'analyse :
\begin{itemize}
    \item \textbf{Le Scout (Externe) :} Cet agent est équipé d'outils (API GitHub, DuckDuckGo, RSS). Sa mission est de capter les signaux faibles du marché et de surveiller les tendances technologiques en temps réel.
    \item \textbf{L'Analyste (Interne) :} Un agent spécialisé dans l'interrogation de la base ChromaDB interne. Sa consigne est stricte : il ne doit jamais utiliser de connaissances externes. S'il ne trouve rien, il doit avouer son ignorance.
    \item \textbf{Le Comparateur :} Il représente l'intelligence analytique du système. Il reçoit les données du Scout et de l'Analyste, et identifie les écarts (gaps) entre la stratégie actuelle de l'entreprise et la réalité mouvante du marché.
\end{itemize}

\section{Les Agents de Production}
Une fois la recherche effectuée, les agents de production prennent le relais :
\begin{itemize}
    \item \textbf{Le Rédacteur :} Formate l'analyse brute en un rapport Markdown professionnel (composé d'un Executive Summary, de risques et d'opportunités).
    \item \textbf{Le Validateur :} Il compare le rapport final généré par le rédacteur aux documents bruts initiaux. Il agit comme un filtre intraitable : il détecte les chiffres inventés ou les affirmations non sourcées.
    \item \textbf{Le Coordinateur :} C'est le routeur conditionnel du graphe. Il gère le point d'entrée et s'assure que le nombre de tentatives (retries) reste limité afin d'éviter les boucles infinies.
\end{itemize}

\chapter{Implémentation Technique Détaillée}

Ce chapitre plonge au cœur du code source (\texttt{pipeline.py}) et décrit, avec un niveau de détail d'ingénierie, les choix techniques et l'architecture logicielle développée.

\section{L'État Partagé : \texttt{VeilleState}}
Dans LangGraph, les agents communiquent via un état partagé. Nous avons défini cet état sous la forme d'un \texttt{TypedDict} en Python :

\begin{lstlisting}[language=Python]
class VeilleState(TypedDict):
    topic: str
    external_context: str
    internal_context: str
    comparison_analysis: str
    final_report: str
    retry_count: int
    logs: Annotated[List[str], operator.add]
    validation: str
    error: Optional[str]
\end{lstlisting}

La gestion de la concurrence est ici primordiale. L'utilisation de \texttt{Annotated[List[str], operator.add]} est une décision d'architecture critique. Elle indique au moteur LangGraph que lorsque plusieurs agents s'exécutent en parallèle (comme le Scout et l'Analyste) et tentent de mettre à jour la clé \texttt{logs} simultanément, le système ne doit pas écraser la donnée (ce qui provoquerait une \texttt{InvalidUpdateError}), mais doit plutôt agir comme un "Reducer" en concaténant les listes. 

\section{Modèles d'IA et Sécurité}
Nous avons utilisé deux modèles open-source distincts via l'API OpenRouter :
\begin{itemize}
    \item \textbf{Modèle Rapide (20B) :} Utilisé pour les requêtes ciblées (Scout, Analyste) nécessitant de la vitesse.
    \item \textbf{Modèle de Raisonnement (120B) :} Réservé au Comparateur et au Rédacteur, dont les tâches exigent de fortes capacités d'analyse conceptuelle.
\end{itemize}

Pour faire face aux contraintes du Cloud (Rate Limits, HTTP 429), nous avons intégré un mécanisme de \textit{Exponential Backoff} natif via le paramètre \texttt{max\_retries=5} de LangChain.

\section{Le Moteur RAG et les Mathématiques Vectorielles}
La base ChromaDB utilise le modèle d'embedding \texttt{all-MiniLM-L6-v2} d'HuggingFace, qui transforme le texte en vecteurs de 384 dimensions. Lors de la récupération des données internes par l'Analyste, nous n'utilisons pas une simple similarité cosinus. Nous implémentons l'algorithme \textbf{MMR (Maximal Marginal Relevance)} :

\begin{lstlisting}[language=Python]
retriever = db.as_retriever(
    search_type="mmr", 
    search_kwargs={"k": 8, "fetch_k": 20, "lambda_mult": 0.7}
)
\end{lstlisting}

L'algorithme MMR récupère d'abord 20 documents très proches mathématiquement, puis extrait les 8 meilleurs en maximisant la diversité (via le multiplicateur \texttt{lambda\_mult}). Cela évite au LLM de lire 8 fois le même paragraphe et élargit considérablement sa vision du contexte de l'entreprise.

\section{Les Outils d'Investigation du Scout}
Le Scout possède un arsenal d'outils décorés par \texttt{@tool}. Chaque fonction contient une docstring précise qui dicte au modèle de langage comment l'utiliser :

\begin{lstlisting}[language=Python]
@tool
def github_trends(query: str):
    """Recherche les depots GitHub les plus populaires sur un sujet donne.
    Utilisez cette fonction pour identifier les frameworks ou 
    projets open-source emergents lies a une technologie."""
    # Requetes HTTP vers l'API GitHub...
\end{lstlisting}
Le Scout parse la docstring et utilise le paradigme de \textit{Function Calling} pour formater ses requêtes JSON afin de déclencher le code Python sous-jacent.

\section{Le Routage Dynamique et le Graphe}
Le système de routage est le "système nerveux" de l'application. 
Le nœud Coordinateur lance la parallélisation (Fan-Out) en retournant explicitement une liste de nœuds à exécuter :
\begin{lstlisting}[language=Python]
def coordinator_router(state):
    # Sécurité de fin
    if "STATUS: PASS" in val or "VALID" in val or state.get("retry_count", 0) >= 2: 
        return "end"
    # Lancement parallele
    return ["scout", "analyst"]
\end{lstlisting}

Le \textbf{Validateur} implémente la boucle anti-hallucination. Son routeur analyse sémantiquement les premiers mots de la sortie du modèle LLM Validateur :
\begin{lstlisting}[language=Python]
def validation_router(state):
    val = state.get("validation", "").upper()
    if "STATUS: FAIL" in val or "ERREURS" in val[:20]: 
        return "writer" # Retour en arriere
    return END
\end{lstlisting}
Si le validateur lève un drapeau d'échec, le flux recule et force le nœud \texttt{writer} à s'exécuter à nouveau en tenant compte des remarques correctives.

\chapter{Développement et Interface Utilisateur (MVC)}

Afin de respecter les principes d'architecture logicielle, le projet suit le modèle MVC (Modèle-Vue-Contrôleur) via une séparation stricte des préoccupations :

\section{Le Backend (Contrôleur et Modèle)}
Le fichier \texttt{pipeline.py} centralise toute la logique IA et la définition du LangGraph. Il expose une fonction \texttt{run\_workflow\_stream} qui agit comme un générateur asynchrone (utilisation de \texttt{yield}). Ce générateur permet d'émettre des événements partiels au fur et à mesure que les nœuds du graphe terminent leur exécution, évitant ainsi de bloquer l'interface pendant l'analyse complète.

\section{Le Frontend (Streamlit)}
Le fichier \texttt{app.py} gère l'interface utilisateur grâce au framework Streamlit. L'UI a été pensée pour être réactive et intuitive. Pendant que le backend travaille, l'utilisateur voit apparaître en temps réel les journaux d'exécution ("Scout en cours", "Analyse Interne terminée", "Vérification anti-hallucination"). Dès que le rapport est généré, il est rendu en Markdown riche directement dans le navigateur.

\section{L'Exportation PDF (ReportLab)}
Le fichier \texttt{pdf\_export.py} prend en charge la transformation du rapport Markdown en un document PDF professionnel. Contrairement à des générateurs basiques, nous utilisons ReportLab. Cela nous permet de gérer finement les styles (polices, marges, espacements) et de convertir les données en objets \texttt{Paragraph} dynamiques pour garantir une mise en page parfaite, même pour les longs paragraphes ou les tableaux complexes.

\chapter{Expérimentations et Défis Techniques}

\section{Scénarios de Validation}
Pour prouver la robustesse de la Cellule de Veille, nous avons soumis le système à différents scénarios complexes :
\begin{itemize}
    \item \textbf{Thèmes complexes et croisés :} "Roadmap IA vs Nouvelles librairies GitHub". L'Analyste a su trouver la roadmap interne, le Scout a ramené les tendances open-source, et le Comparateur a brillamment mis en exergue le retard technologique de l'entreprise.
    \item \textbf{Thèmes hors périmètre :} Lorsque nous avons interrogé le système sur un sujet totalement absent des documents internes de l'entreprise, l'Analyste, strictement bridé par son prompt, a refusé d'inventer des données et a poliment répondu "Information interne non disponible". Cela prouve l'efficacité du verrou anti-hallucination.
\end{itemize}

\section{Résultats Obtenus}
Les performances mesurées sont excellentes :
\begin{itemize}
    \item Génération d'un rapport complet, sourcé et structuré en moins de 10 minutes (souvent quelques dizaines de secondes selon la charge des API).
    \item Fiabilité de 100\% constatée sur l'échantillon testé : les faits avancés par le rapport correspondaient toujours aux sources citées.
    \item L'export PDF automatisé fournit un livrable immédiatement diffusable à la direction.
\end{itemize}

\section{Résolution des Défis Techniques}
Le développement de ce système multi-agents nous a confrontés à deux obstacles majeurs :

\subsection{Défi 1 : La Boucle Infinie d'Hallucination}
\textbf{Problème :} Lors des premières itérations, le LLM jouant le rôle de Validateur (un modèle gratuit de 20B) était trop verbeux. Au lieu de valider simplement le rapport, il rédigeait de longs paragraphes ambigus. Le routeur conditionnel ne comprenait pas la décision et rejetait perpétuellement les bons rapports, créant une boucle infinie d'exécution. \\
\textbf{Solution :} Nous avons appliqué des techniques de Prompt Engineering strictes, obligeant le LLM à répondre EXACTEMENT par \texttt{STATUS: PASS} ou \texttt{STATUS: FAIL}. Le routeur a été modifié pour chercher ces chaînes exactes, rendant la validation déterministe et mettant fin aux faux positifs.

\subsection{Défi 2 : Les Limitations de Quotas API (Erreur HTTP 429)}
\textbf{Problème :} OpenRouter, la passerelle vers nos modèles, bloquait nos requêtes à cause d'une limite de débit (Rate Limit) lorsque les agents Scout et Analyste effectuaient trop de requêtes simultanées. \\
\textbf{Solution :} Outre l'implémentation de \texttt{max\_retries=5} pour l'Exponential Backoff natif, nous avons renforcé le routeur du Validateur pour qu'il intercepte gracieusement les chaînes d'erreurs d'API et termine proprement le graphe (via \texttt{return END}) au lieu de le faire crasher.

\chapter{Conclusion et Perspectives}

\section{Bilan du Projet}
La Cellule SMA (Système Multi-Agents) a rempli toutes ses promesses. Nous avons réussi avec succès la mise en place d'une architecture complexe mêlant LangGraph, bases de données vectorielles (RAG) et Agents autonomes dotés d'outils. L'application Streamlit livrée est fonctionnelle, robuste et immédiatement exploitable par un utilisateur final n'ayant aucune compétence technique. La boucle de validation croisée garantit un niveau de fiabilité inédit par rapport aux approches IA traditionnelles.

\section{Limites Actuelles}
La principale limite de notre solution ne réside pas dans le code, mais dans son architecture Cloud. La dépendance exclusive aux APIs externes (OpenRouter, GitHub, Search) soumet le système à des latences réseau, à des changements de tarification (bien que nous utilisions des modèles gratuits actuellement), et surtout, à des limites de vitesse (Rate Limits) qui empêchent le système d'être utilisé de manière industrielle et massive.

\section{Perspectives d'Évolution et Innovation}
Pour surmonter ces limites et transformer ce projet académique en véritable produit d'entreprise de classe mondiale, plusieurs pistes d'évolutions sont d'ores et déjà identifiées :
\begin{enumerate}
    \item \textbf{Souveraineté totale (On-Premise) :} Héberger les modèles d'IA Open-Source (comme Llama 3) localement sur les serveurs de l'entreprise via des solutions comme \texttt{Ollama}. Cela garantirait un coût nul par requête et une confidentialité absolue des données stratégiques, sans dépendre du Cloud.
    \item \textbf{Connexion aux bases de données structurées :} Améliorer l'agent Analyste en lui donnant des outils capables d'exécuter des requêtes SQL ou de se connecter aux API de Jira et Salesforce, permettant d'inclure des KPI chiffrés et réels dans les rapports de veille.
\end{enumerate}

\end{document}
